\section{Emulating observational effects}

To bridge the gap between the idealized simulation outputs and
realistic survey observations, we employ the RAIL (Redshift Assessment
Infrastructure Layers) software package to emulate observational
effects. RAIL provides a modular framework for injecting realistic
photometric uncertainties, applying survey-specific selection
functions, and simulating the measurement errors characteristic of
modern large-scale imaging surveys. This processing ensures that the
simulated galaxy catalogs reflect the complexities of actual
observations, including magnitude-dependent photometric scatter,
incomplete sky coverage, and the effects of source blending in crowded
fields, thereby providing a more stringent and realistic testbed for
photometric redshift estimation algorithms.


\subsection{Photometric Smearing}

Central to our observational emulation is RAIL's wrapping of the
photometric error module, photErr, which we have extended and wrapped
to account for realistic observing strategies and time-dependent
survey conditions. The standard photErr module provides basic
photometric error modeling based on magnitude-dependent noise
characteristics, but our enhanced version incorporates additional
complexity including spatially varying depth maps. This
wrapper accesses detailed operational simulation outputs that emulate
the expected LSST survey strategy.

Our photErr implementation computes photometric uncertainties by
combining the intrinsic Poisson noise from source photons with
realistic models of sky background, readout noise, and other
systematic contributions. For each simulated galaxy, we use the expected
coadded depth to derive final photometric error estimates. This
approach captures the heterogeneous nature of survey depth across the
footprint, where some regions benefit from numerous high-quality
exposures while others may be observed only during poor
conditions. The resulting photometric uncertainties vary realistically
with position on the sky, band-dependent limiting magnitudes, and
local observing history, providing challenge participants with mock
catalogs whose noise properties more closely match those expected from
the actual survey.


\subsection{Spectroscopic and narrowband photometric redshift
  selection}

RAIL can emulate the selection functions of several different
spectroscopic redshift surveys, including VVDSf02, zCOSMOS,
DEEP2\_LSST, and the DESI BGS, ELG, and LRG samples.

We can also use RAIL to emulate narrowband photometric surveys and
include small amounts of mislabeled reference redshifts.
